% Chapter 1

\chapter{绪论}

\section{研究背景与意义}

近年来，大语言模型（Large Language Models, LLMs）在自然语言处理领域取得了显著进展，并在智能问答、内容生成、代码辅助及智能代理等任务中展现出强大的能力。随着Transformer架构的发展和模型规模的持续扩展，LLMs在复杂推理、多轮交互及跨模态理解等方面不断取得突破，逐渐成为人工智能领域的核心基础技术\cite{touvron2023llama}。同时，模型能力正从单一文本处理向多模态融合与自主决策能力演进，使其在实际应用中的影响力不断扩大。

然而，随着模型能力的增强，其潜在的安全风险也日益凸显。由于训练数据来源复杂且规模庞大，大语言模型不可避免地继承数据中的偏差与不确定性，从而具备生成不安全内容的潜在能力。在缺乏有效约束的情况下，模型可能输出涉及暴力、隐私泄露或误导性信息的内容，对用户与平台安全构成威胁\cite{markov2022holistic,ren2025jailbreakdefensecn,liang2024surveycn}。此外，模型在开放环境中的部署进一步放大了其风险边界，例如在长上下文、多轮交互及工具调用等复杂场景下，模型行为更加难以预测与控制。

近年来，随着大语言模型逐步向智能体（Agent）形态发展，其与外部工具、执行系统及现实任务的耦合程度不断加深。例如，具备自主规划与执行能力的系统使模型不仅能够生成文本，还可能直接影响现实系统决策过程。在此背景下，模型安全问题已从“内容生成风险”演变为“系统级风险”，一旦模型行为失控，可能带来更为严重的连锁影响\cite{hagendorff2026jailbreakagents}。

为降低上述风险，研究界提出了一系列安全对齐（Safety Alignment）方法，如监督微调（Supervised Fine-Tuning, SFT）、基于人类反馈的强化学习（Reinforcement Learning from Human Feedback, RLHF）以及直接偏好优化（Direct Preference Optimization, DPO）等\cite{rafailov2023dpo,hong2024orpo,ethayarajh2024kto}。这些方法通过引入人类价值偏好与安全约束，使模型在生成能力与安全性之间取得一定平衡。然而，现有对齐方法在复杂语境、分布外输入及多样化交互场景中仍存在明显局限，其鲁棒性与泛化能力尚未得到充分保障。

因此，系统研究大语言模型的安全对齐机制及其潜在漏洞具有重要意义。从理论角度来看，有助于深入理解模型对齐过程中的行为约束机制及其失效边界；从实践角度来看，可为风险评估、攻击分析与防御策略设计提供方法基础。在大语言模型持续向更高能力与更广应用场景发展的趋势下，开展相关研究对于提升人工智能系统的安全性与可靠性具有重要价值。

\section{国内外研究现状}

随着大语言模型（Large Language Models, LLMs）的广泛应用，其安全性问题逐渐成为人工智能领域的重要研究方向。近年来，国内外研究围绕安全对齐机制、越狱攻击方法及防御策略等方面展开了系统探索，形成了较为完整的研究体系。现有研究大体可从安全对齐方法、越狱攻击技术以及防御机制三个方面进行归纳。

\subsection{安全对齐方法研究}

为约束大语言模型生成不安全内容，研究界提出了一系列安全对齐方法，主要包括基于数据与基于人类反馈的优化策略。在训练流程上，这些方法通常贯穿预训练、监督微调以及对齐优化等多个阶段。

监督微调（Supervised Fine-Tuning, SFT）通过构建高质量指令数据，使模型学习符合人类期望的行为模式；在此基础上，基于人类反馈的强化学习（Reinforcement Learning from Human Feedback, RLHF）通过引入奖励模型，对模型输出进行偏好引导，从而提升生成结果的安全性与可控性。此外，直接偏好优化（Direct Preference Optimization, DPO）等新方法通过简化优化流程，避免显式奖励模型训练，提高了对齐效率与稳定性\cite{rafailov2023dpo,ethayarajh2024kto}。

除了上述主流方法，研究者还提出多种改进策略，如基于统计优化的对齐方法（RSO、ORPO等），以及结合数据清洗与偏好建模的综合对齐框架\cite{hong2024orpo,markov2022holistic}。这些方法共同推动了大语言模型从“能力导向”向“安全可控”方向发展，使模型在大多数常规场景下能够遵循基本的安全规范。

\subsection{越狱攻击方法研究}

随着对齐技术的发展，研究者逐渐发现，大语言模型在面对特定构造输入时仍可能被诱导绕过安全机制并生成违规内容，这类攻击被称为“越狱攻击”\cite{zou2023universal,chao2024pair,mehrotra2024tap}。相关研究表明，攻击者无需掌握模型内部结构，仅通过设计输入提示即可实现高成功率攻击。

现有越狱攻击方法主要可分为三类：

基于人工构造的越狱方法：通过角色扮演、语义重写或编码转换等方式改写原始问题，使其绕过安全检测\cite{zeng2024johnny,jiang2024artprompt,deng2024multilingual}。

基于优化的对抗攻击方法：利用梯度信息或搜索策略自动生成对抗提示，如GCG（Greedy Coordinate Gradient）方法通过优化输入后缀显著提升攻击成功率\cite{zou2023universal,zhao2024accelerating}。这类方法强调利用模型的概率生成机制实现精确攻击。

基于红队模型的自动化攻击方法：借助辅助模型（Red Teaming LLM）自动生成与优化越狱提示，如GPTFuzzer、PAIR及TAP等方法通过迭代优化显著提升攻击效率与泛化能力\cite{yu2024gptfuzzer,chao2024pair,mehrotra2024tap}。

从更细粒度角度看，越狱攻击还可根据实现机制划分为语言语义攻击、基于优化的对抗攻击以及混合攻击等类型，这些方法从不同层面揭示了模型在语义理解、概率建模及交互过程中的潜在脆弱性。

\subsection{防御方法与安全增强机制研究}

针对越狱攻击带来的风险，研究界提出了多种防御方法，主要可分为模型内部对齐与外部安全增强两类。

在模型内部层面，防御方法主要通过改进训练策略提升模型鲁棒性。例如，通过对抗数据增强，将恶意提示引入训练数据，使模型学习识别潜在攻击；或通过针对性微调提升模型在特定场景下的安全表现\cite{phute2024llmselfdefense,li2023textpurification}。这类方法旨在从参数层面增强模型的安全性。

在模型外部层面，防御策略主要作用于推理阶段，包括输入过滤与输出检测两类机制。输入侧防御通过检测、改写或屏蔽潜在恶意提示，阻止攻击进入模型；输出侧防御则通过审查模型生成结果，确保其符合安全规范。例如，基于规则或辅助模型的输出审查机制已成为实际系统中的常见方案\cite{rebedea2023nemo,markov2022holistic,phute2024llmselfdefense}。

此外，近年来研究逐渐关注全生命周期防御策略，从训练、部署到推理阶段构建统一安全防护体系，以提升模型在复杂应用场景下的整体安全性。

\section{研究问题与本文主要工作}

尽管现有研究在越狱攻击构造、安全对齐训练以及推理阶段防御等方面取得了一定进展，但整体上仍存在若干关键问题亟待系统性解决。现有工作大多从攻击或防御的单一视角出发展开研究，攻防方法之间缺乏统一的分析框架与实验范式，使得不同研究结论难以进行系统对比与综合评估，从而限制了对模型安全机制的整体理解。同时，不同安全对齐方法（如SFT、RLHF、DPO等）在安全性表现上的差异尚缺乏系统性的横向研究，其在目标函数设计、对齐方式及泛化能力方面的内在机制仍不清晰，难以解释模型在复杂语境及分布外输入条件下的对齐失效现象。在攻击侧，现有越狱方法多依赖人工模板或局部优化策略，生成样本在多样性与适应性方面存在局限，难以全面覆盖开放环境中的潜在攻击空间。此外，随着多模态输入与复杂交互场景的逐步引入，跨模态信息与长上下文结构对模型安全性的影响尚未得到充分探讨，使得现有研究在应用层面的解释力仍存在不足。

围绕上述问题，本文以大语言模型安全对齐漏洞为研究核心，从“攻击—评估—防御”的闭环视角出发，对模型安全机制展开系统性分析与方法研究。在机制层面，本文从模型训练与对齐过程入手，对SFT、RLHF及DPO等主流方法在安全约束中的作用机制进行统一分析，重点揭示目标优化之间的潜在冲突、对齐策略的表层性特征以及模型在分布外场景下的泛化局限，从而从理论上解释越狱攻击得以成立的内在原因。在方法层面，针对现有越狱攻击依赖固定模板的问题，设计了一种面向攻击目标的定向越狱提示生成方法，通过结构化建模与自动化生成机制提升攻击样本的多样性与覆盖能力，使攻击过程从经验驱动转向目标驱动。在实验层面，本文构建统一的评测框架，在多模型与多攻击策略条件下对越狱攻击效果进行系统评估，并基于实验结果归纳不同对齐机制下模型的脆弱性特征。在此基础上，进一步结合攻击机理与实验分析结果，从输入层、交互层与输出层三个层面设计分层防御方案，并通过对比实验验证其有效性与适用范围。

基于上述研究，本文在统一分析安全对齐机制的基础上，系统刻画了大语言模型在不同对齐策略下的安全行为特征，并提出了具有更高多样性与可扩展性的越狱攻击生成方法，同时构建了面向攻防一体化的评测框架与分层防御体系，实现了从机制分析到方法设计再到实验验证的闭环研究路径。总体而言，本文从系统视角出发，对大语言模型的安全对齐机制及其越狱攻击问题进行了统一建模与综合分析，推动相关研究由单点方法探索向机制理解与工程化闭环验证转变，为构建更加鲁棒与可信的大模型安全防护体系提供参考。

\section{论文结构安排}

本文围绕大语言模型安全对齐机制及其越狱攻击问题展开研究。全文结构安排如下：

第一章为绪论，介绍研究背景与意义，梳理国内外研究现状，并明确本文的研究问题、研究思路与主要工作。

第二章分析大语言模型安全对齐机制及其潜在漏洞，从大语言模型生成原理与 SFT、RLHF、DPO 等主流对齐方法出发，讨论安全约束在目标冲突、表层对齐和上下文依赖等方面的局限性。

第三章围绕多维越狱攻击框架与实验设计展开研究，在明确越狱攻击威胁模型的基础上，系统组织九类静态攻击维度测试集，并结合自动化红队扩展机制与多轮交互扩展机制说明统一实验流程与评估指标，为后续实验评测提供方法基础。

第四章围绕多维越狱攻击实验评测与结果分析展开，分别从九类静态攻击维度实验、自动化红队补充实验和多轮交互补充实验三个层面评估模型安全性，并结合实验现象对目标冲突、表层对齐与交互过程中的对齐失效机理进行分析。

第五章提出面向越狱攻击的分层防御方案，从输入层、交互层与输出层三个阶段设计协同防御机制，并给出后续用于验证防御效果的实验评估思路。

第六章对全文研究工作进行总结，归纳本文的主要研究内容、创新点与不足，并展望大语言模型安全对齐、越狱攻击评测与鲁棒防御的未来研究方向。
